\section{Rat in a Maze}
\label{sec:rec-rat-in-maze}

\noindent We now begin the chapter's final group: full
\textbf{constraint-satisfaction backtracking}, where the choice at each
step is a move on a grid or board, and a choice already made can block
later choices in ways that must be explicitly tracked and undone.

\problemstatement{Given an $n \times n$ binary maze (cells are either
open, marked $1$, or blocked, marked $0$), a rat starts at the top-left
cell $(0,0)$ and wants to reach the bottom-right cell $(n-1,n-1)$,
moving only through open cells, in any of the four directions (down,
left, right, up), without revisiting a cell already used in the current
path. Find every path from start to destination, each described as a
string of move directions.}

\begin{patternbox}
\emph{``Find all paths through a grid avoiding revisits''} is the
canonical \textbf{grid backtracking} template: at each cell, the choice
set is ``which of up to $4$ neighboring directions to move into next,''
each choice must be validated against the grid's boundaries and
blocked/visited cells \emph{before} being taken, and -- crucially -- a
cell marked as visited while exploring one direction must be
\textbf{unmarked} once that direction's exploration fully returns, so
that a \emph{different} path (reached via backtracking to an earlier
cell) can freely reuse that same cell.
\end{patternbox}

\subsection*{Understanding the problem carefully}

At any cell $(\textit{row},\textit{col})$ during the search, a move into
neighbor $(\textit{newRow},\textit{newCol})$ is valid only if it is
within the grid's bounds, the cell there is open (value $1$), and it has
not already been visited earlier in the \emph{current} path (revisiting
would create a cycle, which is never part of a valid simple path). If
the move is valid, we mark the destination visited, append the
direction's label to the path string, and recurse from there; once that
entire recursive exploration finishes (having found every path reachable
through this move), we must \textbf{unmark} the destination as visited
and remove the direction label, so that a sibling direction's
exploration -- reached by returning to this same cell -- does not
incorrectly believe this cell is permanently blocked.

\begin{ideabox}[Mark Visited Before Recursing, Unmark After Returning]
$\textsc{Solve}(\textit{row}, \textit{col}, \textit{path})$: if
$(\textit{row},\textit{col})$ is the destination: record $\textit{path}$
(and return, or continue if the destination cell itself could lead
further, depending on problem variant). Otherwise, for each of the $4$
directions, in a fixed order: compute the neighbor; if it is in bounds,
open, and not visited: mark it visited; append the direction's label to
$\textit{path}$; recurse into the neighbor; remove the label from
$\textit{path}$ (undo); unmark the neighbor as visited (undo).
\end{ideabox}

\subsection*{Why unmarking after the recursive call returns is essential for correctness}

If a cell, once visited during the exploration of one direction, were
\emph{never} unmarked, it would remain permanently blocked for every
other path that might also legitimately need to pass through it later
-- incorrectly preventing valid paths from being discovered. The
visited-marking only needs to prevent revisiting \emph{within a single,
currently-in-progress path} (to avoid infinite loops/cycles along that
one path); it must not persist across different candidate paths, which
are explored one at a time, each needing a clean slate. Unmarking
immediately after a direction's full recursive exploration returns is
exactly what restores that clean slate for the next direction (or for an
ancestor's next sibling branch) to use.

\subsection*{Algorithm}

\begin{enumerate}
  \item Mark $(0,0)$ visited; call $\textsc{Solve}(0,0,\texttt{""})$.
  \item $\textsc{Solve}(\textit{row}, \textit{col}, \textit{path})$:
  \begin{enumerate}
    \item If $(\textit{row},\textit{col}) = (n-1,n-1)$: record
          $\textit{path}$; return.
    \item For each direction $(dr, dc, \textit{label})$ in $\{(1,0,
          \texttt{D}), (0,-1,\texttt{L}), (0,1,\texttt{R}),
          (-1,0,\texttt{U})\}$:
    \begin{itemize}
      \item $\textit{nr} \gets \textit{row}+dr$, $\textit{nc} \gets
            \textit{col}+dc$.
      \item If $\textit{nr},\textit{nc}$ in bounds, maze cell is open,
            and not visited: mark visited; append $\textit{label}$ to
            $\textit{path}$; recurse $\textsc{Solve}(\textit{nr},
            \textit{nc}, \textit{path})$; remove the appended label
            (undo); unmark visited (undo).
    \end{itemize}
  \end{enumerate}
\end{enumerate}

\subsection*{Dry run}

\dryrun{Take the $2\times2$ maze
$\begin{bmatrix}1&1\\0&1\end{bmatrix}$.}

\begin{itemize}
  \item Start at $(0,0)$, visited$=\{(0,0)\}$,
        $\textit{path}=\texttt{""}$.
  \item Try \texttt{D} (down to $(1,0)$): maze value is $0$ (blocked).
        Invalid, skip.
  \item Try \texttt{L} (left to $(0,-1)$): out of bounds. Skip.
  \item Try \texttt{R} (right to $(0,1)$): open, not visited. Mark
        visited, path$=\texttt{"R"}$. Recurse from $(0,1)$.
  \begin{itemize}
    \item Try \texttt{D} (down to $(1,1)$): open, not visited. Mark
          visited, path$=\texttt{"RD"}$. Recurse from $(1,1)$:
          destination reached! Record \texttt{"RD"}.
    \item Undo: unmark $(1,1)$, path back to $\texttt{"R"}$.
    \item Try \texttt{L} (back to $(0,0)$): visited already, skip.
    \item Try \texttt{U}: out of bounds. Skip.
  \end{itemize}
  \item Undo: unmark $(0,1)$, path back to $\texttt{""}$.
  \item Try \texttt{U} (up to $(-1,0)$): out of bounds. Skip.
\end{itemize}

The only path found is \texttt{"RD"} (right, then down).

\subsection*{C++ implementation}

\begin{lstlisting}
#include <bits/stdc++.h>
using namespace std;

void solve(int row, int col, vector<vector<int>>& maze, int n,
           vector<vector<bool>>& visited, string& path,
           vector<string>& result) {
    if (row == n - 1 && col == n - 1) {
        result.push_back(path);
        return;
    }

    int dr[] = {1, 0, 0, -1};
    int dc[] = {0, -1, 1, 0};
    char label[] = {'D', 'L', 'R', 'U'};

    for (int d = 0; d < 4; d++) {
        int nr = row + dr[d], nc = col + dc[d];
        if (nr >= 0 && nr < n && nc >= 0 && nc < n &&
            maze[nr][nc] == 1 && !visited[nr][nc]) {
            visited[nr][nc] = true;
            path.push_back(label[d]);

            solve(nr, nc, maze, n, visited, path, result);

            path.pop_back();
            visited[nr][nc] = false;
        }
    }
}

vector<string> ratInMaze(vector<vector<int>>& maze) {
    int n = maze.size();
    vector<string> result;
    if (maze[0][0] == 0) return result;

    vector<vector<bool>> visited(n, vector<bool>(n, false));
    visited[0][0] = true;
    string path;
    solve(0, 0, maze, n, visited, path, result);
    return result;
}

int main() {
    vector<vector<int>> maze = {{1,1},{0,1}};
    for (string& p : ratInMaze(maze)) cout << p << " ";
    cout << endl;
    return 0;
}
\end{lstlisting}

\begin{complexitybox}
\textbf{Time Complexity.} In the worst case (an open grid with no
blocked cells), the number of simple paths can be exponential in $n^2$
(roughly $O(4^{n^2})$ branching, though heavily pruned by the
no-revisit rule in practice); a tight closed-form bound is not generally
useful here -- what matters is that each cell visit does $O(1)$ work
beyond the recursive calls themselves.
\textbf{Space Complexity.} $O(n^2)$ for the visited grid, plus
$O(n^2)$ for the maximum recursion depth (a path can visit at most every
cell once).
\end{complexitybox}

\begin{takeawaybox}
``Mark visited before recursing, unmark after returning'' is the
universal template for any backtracking search over a grid or graph that
must avoid revisiting nodes within a single path -- it is the direct
generalization of the choose/recurse/un-choose rhythm from
Section~\ref{sec:rec-power-set}, applied to spatial moves instead of
array-element inclusion. The same template, with different validity
checks, drives N-Queens, Sudoku, and Word Search in the sections that
follow.
\end{takeawaybox}
